\section{Перспективы}
\label{sec:outlook}

\textbf{Следующая версия.} Четыре замера переведут главные оценки работы в измерения. Первый: многоходовая сессия
Qwen3.8-27B между сервером на CUDA и клиентом на Metal с контекстом до 256 тысяч токенов и локальными ходами между
оффлоадами. Второй: качество одинаковых и разнородных квантов на агентских задачах. Третий: полная задержка вызова с
очередью, подъёмом копии, передачей и вставкой при нескольких клиентах одновременно. Четвёртый: сравнение с
современным локальным движком и с сервером подходящего размера при одинаковых требованиях к качеству и скорости.
Кроме того, следующая версия добавит замеры на других клиентах, начиная с M5 Max 64~ГБ, и статистику агентских
сессий других участников.

\textbf{Пул простаивающих мощностей.} Своя карта дешевле аренды, только когда она загружена, а одна команда в одном часовом поясе использует её в
рабочие часы, около четверти календарного времени (8 часов в 21 рабочий день). Если сервер гибрида может сдавать простой в общий пул и
мягко забирать его обратно, доход от пула может снизить стоимость своей карты. Насколько, зависит от спроса, цены и
доли времени, которую можно продать без вытеснения рабочих сессий. Пул похож на OpenRouter, но
для узкого класса моделей, где поставщики взаимозаменяемы. У такого пула есть нерешённые задачи: состояние сессии
живёт на узле, поставщик видит токены чужой команды, а целостность его работы надо проверять. Это отдельное
направление исследований.

\textbf{Проверяемый prefill как отрасль.} Если клиенты научатся проверять память модели (раздел~\ref{sec:security}),
поставщики смогут продавать проверяемый prefill тем, кому важна безопасность.

\textbf{Каждая сторона на своём лучшем движке.} Сервер может считать на самых быстрых для своего железа форматах,
например 4-битных NVFP4 на NVIDIA, а клиент генерировать на MLX (ждёт проверки). Для этого нужен мост между форматами кэша двух
движков и проверка эквивалентности для каждой пары.

\textbf{Сжатие дельты.} Передача рекуррентного состояния в 16 битах вдвое уменьшает постоянную часть дельты и
расширяет долю вызовов, где оффлоад выгоден (таблица~\ref{tab:offload}); её нужно проверить на эквивалентность
генерации.
